% !TEX program = xelatex
% 格林函数：恒等式、三类边界、一维例子
\documentclass[fontset=windows, chinese, 10pt, xcolor=dvipsnames]{beamer}
% Shared preamble for the chapter 10 standalone lectures.
% Each lecture sets \title and \subtitle after this file is read.

\usepackage{amsmath,amssymb,mathtools}
\usepackage{xeCJK}
\usepackage{zhnumber}
\usepackage{tikz}
\usepackage{tikz-3dplot}
\usetikzlibrary{arrows.meta,calc,positioning}

\usetheme{metropolis}
\usepackage{beamercolorthemeHistNanjing}

\newcommand{\red}[1]{\textcolor{red}{#1}}
\newcommand{\blue}[1]{\textcolor{blue}{#1}}
\newcommand{\green}[1]{\textcolor{green}{#1}}
\newcommand{\gray}[1]{\textcolor{gray}{#1}}
\newcommand{\orange}[1]{\textcolor{orange}{#1}}

\newcommand{\sssub}[1]{ {\mathrm{#1} }}  % script size subscript\\\
\newcommand{\funcd}[2]{ \frac{d{#1}}{d{#2}}} %function full derivative
\newcommand{\funcpd}[2]{ \frac{\partial{#1}}{\partial{#2}}} %function partial derivative
\newcommand{\fnald}[2] { \frac{\delta {#1}}{\delta {#2}}} % functional derivative

\renewcommand{\Re}{\operatorname{Re}}
\renewcommand{\Im}{\operatorname{Im}}
\def\ext{\sssub{ext}}
\def\x{\sssub{x}}
\def\xc{\sssub{xc}}
\def\H{\sssub{H}}
\def\Hx{\sssub{Hx}}
\def\Hxc{\sssub{Hxc}}
\def\c{\sssub{c}}
\def\s{\sssub{s}}

% \def\ext{\sssub{ext}}
% \def\ext{\sssub{ext}}
% \def\ext{\sssub{ext}}
% \def\ext{\sssub{ext}}
% \defsssub{ext}


% definitions of bold faced symbols
\def\brpppp{{\mathbf{r}^{\prime\prime\prime\prime}}}
\def\brppp{{\mathbf{r}^{\prime\prime\prime}}}
\def\brpp{{\mathbf{r}^{\prime\prime}}}
\def\brp{{\mathbf{r}^{\prime}}}
\def\bzp{{\mathbf{z}^{\prime}}}
\def\bxp{{\mathbf{x}^{\prime}}}
\def\tp{{{t}^{\prime}}}
\def\tpp{{{t}^{\prime\prime}}}
\def\tppp{{{t}^{\prime\prime\prime}}}

\def\tbr{{\tilde{\mathbf{r}}}}
\def\bk{{\mathbf{k}}}
\def\br{{\mathbf{r}}}
\def\bp{{\mathbf{p}}}
\def\bq{{\mathbf{q}}}
\def\bv{{\mathbf{v}}}
\def\bz{{\mathbf{z}}}
\def\mbf{{\mathbf{f}}}

\def\bx{{\mathbf{x}}}
\def\bR{{\mathbf{R}}}
\def\bM{{\mathbf{M}}}
\def\bP{{\mathbf{P}}}
\def\bT{{\mathbf{T}}}
\def\bK{{\mathbf{K}}}
\def\bA{{\mathbf{A}}}
\def\bB{{\mathbf{B}}}
\def\bD{{\mathbf{D}}}
\def\bE{{\mathbf{E}}}
\def\bF{{\mathbf{F}}}
\def\bG{{\mathbf{G}}}
\def\bH{{\mathbf{H}}}
\def\bI{{\mathbf{I}}}
\def\bJ{{\mathbf{J}}}
\def\bP{{\mathbf{P}}}
\def\bS{{\mathbf{S}}}
\def\bV{{\mathbf{V}}}
\def\bX{{\mathbf{X}}}
\def\bY{{\mathbf{Y}}}
\def\dd{{\mathop{}\!\mathrm{d}}}
\def\dr{{\!\!\dd\br\,}}
\def\dx{{\!\!\dd\bx\,}}
\def\dy{{\!\!\dd\by\,}}
\def\dk{{\!\!\dd\bk\,}}
\def\dG{{\!\!\dd\bG\,}}


\def\half{{\frac{1}{2}}}
\def\tr{{\mathrm{tr}}}


\def\Arg{{\mathrm{Arg}\;}}
\def\arg{{\mathrm{arg}\;}}
\def\Res{{\mathrm{Res}\;}}

% unit-vector basis (bold under XeLaTeX; \mathbf can look non-bold)
\newcommand{\ue}{\boldsymbol{e}}
\newcommand{\un}{\hat{\boldsymbol{n}}}


\author{罗凯}
\institute{南京理工大学物理学院}
\date{\zhdigits{\the\year},\zhganzhinian{\the\year}年秋}

\usebackgroundtemplate{%
\begin{tikzpicture}[remember picture,overlay]
  \node[opacity=0.05, at=(current page.center)] {
    \includegraphics[width=0.6\paperwidth]{njust-logo.jpeg}
  };
\end{tikzpicture}}


\title{格林函数法}
\subtitle{第一讲：格林恒等式}

\begin{document}

\begin{frame}
  \titlepage
\end{frame}

\begin{frame}{本讲提纲}
  分离变量把定解问题拆成本征模态。格林函数走另一条路：先求单位点源在同样边界下的场，再线性叠加。
  \begin{enumerate}
    \item 方程写成 $\nabla^2 u=f$，基本解取负号；
    \item 格林第二恒等式，挖掉点源处的小球；
    \item 三类边界各留下一个积分公式；
    \item 区间 $(0,l)$ 上把公式算到底。
  \end{enumerate}
\end{frame}

\section{恒等式}

\begin{frame}{符号约定}
  这一章的泊松方程写成
  \begin{equation*}
    \boxed{\displaystyle \nabla^2 u=f.}
  \end{equation*}
  自由空间的基本解因此是负的，
  \[
    \nabla^2\Bigl(-\frac{1}{4\pi R}\Bigr)=\delta^3(\mathbf{r}).
  \]
  若方程改成 $-\nabla^2 u=\rho$，格林函数整体差一个负号，下面每个公式都要改号。
\end{frame}

\begin{frame}{格林第一恒等式}
  对 $u\nabla v$ 用散度定理，
  \[
    \iint_{\Sigma} u\,\nabla v\cdot\mathrm{d}\mathbf{S}
    =\iiint_{T}\bigl(u\nabla^2 v+\nabla u\cdot\nabla v\bigr)\,\mathrm{d}V.
  \]
  $\partial/\partial n$ 取区域 $T$ 的外法向。$u$ 和 $v$ 在 $T$ 上要有连续的二阶导数，边界分段光滑。
\end{frame}

\begin{frame}{格林第二恒等式}
  把 $u$ 和 $v$ 对调再相减，梯度内积消去，
  \begin{equation*}
    \boxed{\displaystyle
    \iint_{\Sigma}\Bigl(
      u\frac{\partial v}{\partial n}-v\frac{\partial u}{\partial n}
    \Bigr)\,\mathrm{d}S
    =\iiint_{T}\bigl(u\nabla^2 v-v\nabla^2 u\bigr)\,\mathrm{d}V.}
  \end{equation*}
  点源在 $\mathbf{r}_0$ 处不可导，不能把 $u$ 和 $v$ 直接代进去。
\end{frame}

\section{点源}

\begin{frame}{三类边界}
  边界上给定
  \[
    \Bigl[\alpha\frac{\partial u}{\partial n}+\beta u\Bigr]_{\Sigma}=\varphi.
  \]
  $\alpha=0$ 是狄利克雷，$\beta=0$ 是诺伊曼，两者都不为零是罗宾。
  \vspace{0.4em}
  单位点源的场满足 $\nabla^2 v=\delta(\mathbf{r}-\mathbf{r}_0)$。它在源点发散，先挖掉以 $\mathbf{r}_0$ 为心、半径 $\varepsilon$ 的小球。
\end{frame}

\begin{frame}{挖去小球}
  剩下的区域里 $v$ 调和，$\nabla^2 u=f$。第二恒等式写成
  \begin{align*}
    &\iint_{\Sigma}\Bigl(v\partial_n u-u\partial_n v\Bigr)\,\mathrm{d}S
    +\iint_{\Sigma_{\varepsilon}}\Bigl(v\partial_n u-u\partial_n v\Bigr)\,\mathrm{d}S\\
    &\qquad=\iiint_{T\setminus K_{\varepsilon}} v\,f\,\mathrm{d}V.
  \end{align*}
  小球的外法向指向球心，$\partial/\partial n=-\partial/\partial\varepsilon$。附近 $v\sim -1/(4\pi|\mathbf{r}-\mathbf{r}_0|)$。
\end{frame}

\begin{frame}{小球上的极限}
  $v$ 是 $O(1/\varepsilon)$，面积是 $4\pi\varepsilon^2$，所以
  \[
    \iint_{\Sigma_{\varepsilon}} v\frac{\partial u}{\partial n}\,\mathrm{d}S\to 0.
  \]
  $\partial v/\partial\varepsilon\sim 1/(4\pi\varepsilon^2)$，法向反向，
  \[
    \iint_{\Sigma_{\varepsilon}} u\frac{\partial v}{\partial n}\,\mathrm{d}S\to -u(\mathbf{r}_0).
  \]
  令 $\varepsilon\to 0$，
  \begin{equation*}
    u(\mathbf{r}_0)
    =\iiint_T v\,f\,\mathrm{d}V
    -\iint_{\Sigma}\Bigl(
      v\partial_n u-u\partial_n v
    \Bigr)\,\mathrm{d}S.
  \end{equation*}
\end{frame}

\begin{frame}{未知的边界值要消掉}
  上面的公式同时用到边界上的 $u$ 和 $\partial u/\partial n$。定解问题只给它们的一个线性组合。
  \vspace{0.4em}
  把 $v$ 选成满足同样类型的齐次边界条件，未知的那一组边界值就不再出现。这样选出的 $v$ 记作格林函数 $G$。
\end{frame}

\section{三类公式}

\begin{frame}{狄利克雷}
  要求 $G|_{\Sigma}=0$。表示公式变成
  \begin{equation*}
    \boxed{\displaystyle
    u(\mathbf{r}_0)
    =\iiint_T G\,f\,\mathrm{d}V
    +\iint_{\Sigma}\phi\,\frac{\partial G}{\partial n}\,\mathrm{d}S,}
  \end{equation*}
  其中 $\phi=u|_{\Sigma}$。$f=0$ 时只剩面积分：调和函数由边界值经 $\partial G/\partial n$ 再生。
\end{frame}

\begin{frame}{罗宾}
  $G$ 满足同一系数的齐次条件，边界上 $\partial G/\partial n=-(\beta/\alpha)G$。代入后，含未知 $u$ 的面积分两两抵消，
  \[
    u(\mathbf{r}_0)
    =\iiint_T G f\,\mathrm{d}V
    -\frac{1}{\alpha}\iint_{\Sigma} G\,\psi\,\mathrm{d}S.
  \]
  $\psi$ 是边界上给定的那个线性组合。
\end{frame}

\begin{frame}{诺伊曼要相容}
  不能同时要求 $\nabla^2 G=\delta$ 和 $\partial G/\partial n|_{\Sigma}=0$。对方程积分，左边是边界上法向导数的积分，等于零，右边是 $1$。
  \vspace{0.4em}
  三维里改成
  \[
    \nabla^2 G=\delta(\mathbf{r}-\mathbf{r}_0)-\frac{1}{V_T},
    \qquad
    \frac{\partial G}{\partial n}\Big|_{\Sigma}=0.
  \]
  $V_T$ 是区域体积。均匀汇把点源的总量抵消。二维用面积代替体积。原问题也要满足 $\iiint f=\iint\partial u/\partial n$。
\end{frame}

\begin{frame}{互易}
  \begin{equation*}
    \boxed{\displaystyle G(\mathbf{r},\mathbf{r}_0)=G(\mathbf{r}_0,\mathbf{r}).}
  \end{equation*}
  两个源点都挖成小球，对 $G(\cdot,\mathbf{r}_1)$ 和 $G(\cdot,\mathbf{r}_2)$ 用第二恒等式。体积内被积函数是零，外边界上齐次条件使面积分为零，两个小球的极限留下 $G(\mathbf{r}_1,\mathbf{r}_2)-G(\mathbf{r}_2,\mathbf{r}_1)=0$。
  \vspace{0.35em}
  观察点和源点可以换写。常用形式是对 $\mathbf{r}_0$ 积分，在 $\mathbf{r}$ 处读出 $u$。
\end{frame}

\section{一维}

\begin{frame}{区间上的格林函数}
  $0<x<l$，$u''=f$，$u(0)=u(l)=0$。两侧是直线，
  \[
    G=
    \begin{cases}
      A x, & 0\leq x\leq x_0,\\
      C(x-l), & x_0\leq x\leq l.
    \end{cases}
  \]
  连续给出 $A x_0=C(x_0-l)$。越过 $x_0$ 积分，$G'(x_0+)-G'(x_0-)=1$，即 $C-A=1$。
\end{frame}

\begin{frame}{解出来}
  \begin{equation*}
    \boxed{\displaystyle
    G(x,x_0)=
    \begin{cases}
      (x_0-l)x/l, & x\leq x_0,\\
      x_0(x-l)/l, & x\geq x_0.
    \end{cases}}
  \end{equation*}
  它对 $x$ 和 $x_0$ 对称，区间内部为负，和基本解 $-1/(4\pi R)$ 的符号一致。
  \vspace{0.35em}
  取 $f=-1$，则 $u=\int_0^l G(x,x_0)(-1)\,\mathrm{d}x_0=x(l-x)/2$。直接积分 $u''=-1$ 得到同一条抛物线。
\end{frame}

\begin{frame}{这一讲留下的公式}
  方程是 $\nabla^2 u=f$，基本解是 $-1/(4\pi R)$。
  \vspace{0.35em}
  狄利克雷格林函数在边界上为零，
  \[
    u(\mathbf{r}_0)
    =\iiint G f\,\mathrm{d}V
    +\iint\phi\,\partial_n G\,\mathrm{d}S.
  \]
  诺伊曼问题要先把源的总量用均匀汇抵消。一维区间上，$G$ 是折线，斜率在源点跳 $1$。
\end{frame}

\end{document}
